\documentclass[10pt,a4paper]{amsart}
\usepackage[margin=19mm]{geometry}
\usepackage{amsmath,amssymb,mathtools}
\usepackage[scr=rsfs,cal=euler]{mathalfa}
\usepackage{xfrac}
\usepackage[unicode,hidelinks,bookmarks=false]{hyperref}
\newcommand{\ie}{i.e.\ }
\newcommand{\cf}{cf.\ }
\newcommand{\resp}{respectively\ }
\newcommand{\Subsection}{\subsection}
\providecommand{\nobreakdash}{\nobreak\hbox{-}}
\setlength{\emergencystretch}{2em}
\multlinegap0pt
\usepackage{amssymb}
\newtheorem{lemma}{Lemma}
\title{{\bf Diagonal Multi-omics Integration of Heterogeneous Datasets}}
\author{Maksim \, V.~Kukushkin, Mikhail S. Arbatskiy, Dmitriy E. Balandin, }
\date{Source: arXiv:2608.16968v2 (CC BY 4.0)}
\begin{document}
\maketitle

\noindent\emph{Source and licence.} {\bf Diagonal Multi-omics Integration of Heterogeneous Datasets}. Version arXiv:2608.16968v2 (CC BY 4.0), distributed by arXiv under the Creative Commons Attribution 4.0 International licence, http://creativecommons.org/licenses/by/4.0/, as recorded in the arXiv licence metadata for this exact version and shown on the abstract page https://arxiv.org/abs/2608.16968v2. Full licence notice: \texttt{licenses/arxiv-2608.16968-CC-BY-4.0.txt}.\par\smallskip

\noindent\textit{Editorial scope.} Counted unit: the article's lemma L3 (source0.tex lines 536--540). The article's complete original proof of it is delivered with it (source0.tex lines 541--623, one \texttt{proof} environment of 83 lines). No mathematics is written, repaired or re-derived: the statement and the proof are the article's own lines, in the article's own notation, and no retained line is substituted or re-flowed.  Retained uncounted as the source-local context the proof needs: lines 261--263; lines 412--414; lines 416--418; lines 468--473.  Nothing was omitted from any retained span, so the package carries no omission record.  No retained line contains a citation, so the wrapper carries no bibliography and no bibliography entry.  No retained line contains a figure environment, an image-inclusion command or drawing code, so no image is delivered and the package declares no asset dependency.  Editorial formatting only, in addition: an AMS \texttt{amsart} preamble that restates the article's own declarations and macros used by the retained lines, namely the environments \texttt{lemma} on separate counters, together with the standard packages those lines need.  The retained environments print in this document as Lemma 1 in order of appearance, whereas the article prints the same items with the numbering of its own full document.  The complete native source of the article as distributed in the arXiv e-print for this exact version is bundled unchanged as \texttt{sources/207865/source0.tex}.
\begin{equation}\label{2}
X^{\ast}T_{\zeta}X= I_{m},\;X\in \mathbb{C}^{n\times m},\,m<n,
\end{equation}

\begin{equation}\label{7}
\Psi(\theta)\mathbf{e}_{ j } =\mu_{j} T_{\zeta}\mathbf{e}_{ j },\;(T_{\zeta}\mathbf{e} _{ j },\mathbf{e} _{ k })_{\mathbb{C}^{n}}=\delta_{jk},\;j,k=1,2,\dots,\xi,
\end{equation}

\begin{equation}\label{8}
T(\theta)\mathbf{v}_{j} =\mu_{j}\mathbf{v}_{j},\;(\mathbf{v}_{j},\mathbf{v}_{k})_{\mathbb{C}^{n}}=\delta_{jk},\,\mathbf{v}_{j}=\sqrt{T_{\zeta}}\mathbf{e}_{j},\;j,k=1,2,\dots,\xi,
\end{equation}

\begin{lemma}\label{L2}
The matrix functions $E_{p}(\theta)$ are analytic on the real axis. Moreover, there exists a unique  solution to the following system
\begin{equation}\label{9}
\mathrm{arg}\, \psi\{E_{p}(\theta)\}=\theta,\;\theta\in \mathbb{J},\;p=1,2,\dots,C_{n-1}^{m}.
\end{equation}
\end{lemma}

\begin{lemma}\label{L3} The following relation holds
\begin{equation}\label{10}
 \max\limits_{X^{\ast}T_{\zeta}X= I_{m}} \sum\limits_{j=1}^{m}(\Psi(\theta)\mathbf{x}_{j},\mathbf{x}_{j})_{\mathbb{C}^{n}}=\sum\limits_{s=n-m }^{n-1}\mu_{s}(\theta),\;\theta\in (0,\pi/2).
\end{equation}
\end{lemma}

\begin{proof} For simplicity of notation,  we assume that the eigenvalues  are arranged  in decreasing order of their absolute values.
 Applying Lemma \ref{L2}, we can choose  $T_{\zeta}$-orthonormal analytic eigenvectors $\{\mathbf{e}_{j}\}^{\xi}_{1}.$
Using the property of the selfadjoint operator $T(\theta)$ and taking into account the equivalence between the eigenvalue problems \eqref{7} and  \eqref{8}, we can easily prove
$$
\mathbb{C}^{n}=\mathrm{span}\{\mathbf{1}\}\oplus_{\,T_{\zeta}}\mathrm{span}\{\mathbf{e}_{1},\mathbf{e}_{2},\ldots,\mathbf{e}_{n-1}\}.
$$
Thus,  the following representation holds
$$
 \sum\limits_{j=1}^{m}( \Psi(\theta) \mathbf{x} _{j}  ,     \mathbf{x} _{j}   )_{\mathbb{C}^{n}}=  \sum\limits_{j=1}^{m}\left( \sum\limits_{s=1}^{\xi}c_{sj } \mu_{s}
          T_{\zeta}\mathbf{e}_{s},\sum\limits_{s=1}^{\xi} c_{sj } \mathbf{e}_{s} \right)_{\! \mathbb{C}^{n}}=
          \sum\limits_{s=1}^{\xi}\mu_{s} \sum\limits_{j=1}^{m}  |c_{sj }|^{2},
$$
where we have used the decomposition
$$
\mathbf{x} _{j}= \sum\limits_{s=1}^{n-1} c_{sj }\mathbf{e}_{s}+c_{nj} \mathbf{1},
$$
the values $\mu_{s},\,c_{sj },$ and the vectors $\mathbf{e}_{s}$ depend on the parameter $\theta,$ according to the   agreement, we have $\mu_{1}>\mu_{2}>...>\mu_{\xi}.$ It is clear that the eigenvalues of the problem  \eqref{7} are positive numbers. Using the connection condition \eqref{2}, we get
\begin{equation*}
 \delta_{jk}=(  T_{\zeta} \mathbf{x} _{j} ,\mathbf{x} _{k}   )_{\mathbb{C}^{n}}=
 \left(\sum\limits_{s=1}^{n}c_{sj }T_{\zeta}\mathbf{e}_{s},\sum\limits_{s=1}^{n} c_{sk }\mathbf{e}_{s} \right)_{\!\!\!\mathbb{C}^{n}}=
 \sum\limits_{s=1}^{\xi} c _{sj }\overline{c} _{sk }+c_{nj} \bar{c}_{nk}\,\mathrm{tr}(T_{\zeta}),
\end{equation*}
where $\mathbf{e}_{n}=\mathbf{1}.$
Applying this formula, rewriting the left-hand side of \eqref{10}, we come to the related  problem
$$
\max\limits \sum\limits_{s=1}^{\xi}\mu _{s} \sum\limits_{j=1}^{m}  |c_{sj }|^{2},
\;\;\sum\limits_{j=1}^{m}\sum\limits_{s=1}^{\xi} |c _{sj }|^{2}+\mathrm{tr}(T_{\zeta})\sum\limits_{j=1}^{m}|c _{nj }|^{2}=m,
$$
Denote
$$
u_{0}:=\mathrm{tr}(T_{\zeta})\sum\limits_{j=1}^{m}|c _{nj }|^{2},\;u_{s}:=\sum\limits_{j=1}^{m}  |c_{sj }|^{2},\;s=1,2,\ldots,\xi.
$$
Then  the expression corresponding to the  main problem can be rewritten in the form
$$
d_{\xi}(\mathbf{u}) :=\sum\limits_{s=1}^{\xi}\mu _{s}u_{s}  ,\;\;\sum\limits_{s=1}^{\xi}u_{s}+u_{0}=m,\;\mathbf{u}=(u_{0},u_{1},\ldots,u_{\xi}).
$$
Substituting, we get
$$
h_{\xi-1}(\mathbf{u}):=\sum\limits_{s=1}^{\xi-1}u_{s}\mu _{s}+\mu _{\xi}\left(m-\sum\limits_{s=1}^{\xi-1}u_{s}-u_{0}\right)=
\sum\limits_{s=1}^{\xi-1}u_{s} \left(\mu _{s}-\mu_{\xi}\right)-u_{0}\mu _{\xi}+m \mu _{\xi},
$$
$$
\sum\limits_{s=0}^{\xi-1}u_{s}\leq m.
$$
Therefore, the component of the gradient $\nabla h_{\xi-1}$ corresponding to   $u_{0}$ is   negative, whereas the  other  coordinates are positive  $ \mu_s - \mu_{\xi} > 0.$ Combining this fact with the geometric properties of the hyperplane, we immediately conclude that at the maximum point the non-negative parameter $u_{0}$ must take  its lower bound, i.e., $u_{0} = 0.$ The latter eliminates the kernel contribution. Thus, we  have
$$
\max \limits_{\mathbf{u}\in \mathbb{R}^{\xi-1}_{m}}  h_{\xi-1}(\mathbf{u})= \max\limits_{\mathbf{u}\in \mathbb{T}^{\xi-1}_{m}}\sum\limits_{s=1}^{\xi-1}u_{s}\mu_{s}=  \max\limits_{\mathbf{u}\in \mathbb{T}^{\xi-1}_{m}}  d_{\xi-1}(\mathbf{u}),
$$
where
$$
\mathbb{R}^{n}_{m}:=\left\{\mathbf{x}\in \mathbb{R}^{n}:\, \sum\limits_{s=1}^{n}x_{s}\leq m \right\},\;\mathbb{T}^{n}_{m}:=\left\{\mathbf{x}\in \mathbb{R}^{n}:\, \sum\limits_{s=1}^{n}x_{s}= m \right\},\;n,m\in \mathbb{N}.
$$
Apparently, in finding the maximum point, we have some freedom in choosing the coefficients, thus at this step of the proof, we   assume    $c_{\xi j}=0,\,j=1,2,\ldots,m.$
Following  the same reasoning, we get
\begin{equation*}
h_{\xi-2}(\mathbf{u}):=
\sum\limits_{s=1}^{\xi-2}u_{s} \left(\mu _{s}-\mu _{\xi-1}\right)+m \mu _{\xi-1},\;\sum\limits_{s=1}^{\xi-2}u_{s}\leq m.
\end{equation*}
Therefore, the constant vector $\nabla h_{\xi-2}$ has positive coordinates. Analogously to the above, we get
$$
\max\limits_{\mathbf{u}\in \mathbb{R}^{\xi-2}_{m}}  h_{\xi-2}(\mathbf{u})=   \max\limits_{\mathbf{u}\in \mathbb{T}^{\xi-2}_{m}}  d_{\xi-2}(\mathbf{u}),
$$
where the latter relation is obtained due to the choice  $c_{\xi-1j}=c_{\xi j}=0,\,j=1,2,\ldots,m.$
Following the same reasoning, we come to the problem
$$
\max\limits_{X^{\ast}T_{\zeta}X= I} \sum\limits_{j=1}^{m}(\Psi(\theta)\mathbf{x}_{j},\mathbf{x}_{j})_{\mathbb{C}^{n}}=\max\limits_{ C ^{\ast} C =I_{m}} \sum\limits_{s=1}^{m}\mu _{s} \sum\limits_{j=1}^{m}  |c_{sj }|^{2},
$$
where we assume  $c_{m+1j}=c_{m+2j}=\ldots=c_{\xi j}=0,\,j=1,2,\ldots,m.$
Therefore, we can restrict the problem considering the   relation
\begin{equation}\label{11}
      \sum\limits_{s=1}^{m}\mu_{s} \sum\limits_{j=1}^{m}  |c_{sj }|^{2}, \;
\sum\limits_{s=1}^{m} c _{sj }\overline{c }_{sk }  =\delta_{kj},\,j,k=1,2,\ldots,m.
\end{equation}
 In  matrix form, we have
$
 C ^{\ast} C =I_{m},\;  \Rightarrow  C C ^{\ast}=I_{m},\;C:=\{c_{sj}\}\in \mathbb{C}^{m\times m},
$
i.e.,
$$
\sum\limits_{j=1}^{m} c _{sj  }\overline{c }_{kj } =\delta_{sk},\,s,k=1,2,\ldots,m.
$$
Substituting   into \eqref{11} and    arranging   the eigenvalues in increasing order of their absolute values, we obtain \eqref{10}. The proof is complete.
\end{proof}

\end{document}
